\section{Serre's criterion for normality}
\label{section-serre-criterion}

\noindent
We introduce the following properties of Noetherian rings.

\begin{definition}
\label{definition-conditions}
Let $R$ be a Noetherian ring.
Let $k \geq 0$ be an integer.
\begin{enumerate}
\item We say $R$ has property {\it $(R_k)$} if for every prime $\mathfrak p$
of height $\leq k$ the local ring $R_{\mathfrak p}$ is regular.
We also say that $R$ is {\it regular in codimension $\leq k$}.
\item We say $R$ has property {\it $(S_k)$} if for every prime $\mathfrak p$
the local ring $R_{\mathfrak p}$ has depth at least
$\min\{k, \dim(R_{\mathfrak p})\}$.
\item Let $M$ be a finite $R$-module. We say $M$ has property $(S_k)$
if for every prime $\mathfrak p$ the module
$M_{\mathfrak p}$ has depth at least
$\min\{k, \dim(\text{Supp}(M_{\mathfrak p}))\}$.
\end{enumerate}
\end{definition}

\noindent
Any Noetherian ring has property $(S_0)$ and so does any finite module
over it. Our convention that the depth of the zero module is $\infty$
(see Section \ref{section-depth}) and the dimension of the empty set is
$-\infty$ (see Topology, Section \ref{topology-section-krull-dimension})
guarantees that the zero module has property $(S_k)$ for all $k$.

\begin{lemma}
\label{lemma-criterion-no-embedded-primes}
Let $R$ be a Noetherian ring.
Let $M$ be a finite $R$-module.
The following are equivalent:
\begin{enumerate}
\item $M$ has no embedded associated prime, and
\item $M$ has property $(S_1)$.
\end{enumerate}
\end{lemma}

\begin{proof}
Let $\mathfrak p$ be an embedded associated prime of $M$.
Then there exists another associated prime $\mathfrak q$ of $M$
such that $\mathfrak p \supset \mathfrak q$. In particular this
implies that $\dim(\text{Supp}(M_{\mathfrak p})) \geq 1$ (since $\mathfrak q$
is in the support as well). On the other hand $\mathfrak pR_{\mathfrak p}$
is associated to $M_{\mathfrak p}$
(Lemma \ref{lemma-associated-primes-localize}) and hence
$\text{depth}(M_{\mathfrak p}) = 0$
(see Lemma \ref{lemma-ideal-nonzerodivisor}).
In other words $(S_1)$ does not hold.
Conversely, if $(S_1)$ does not hold then there exists a prime
$\mathfrak p$ such that $\dim(\text{Supp}(M_{\mathfrak p})) \geq 1$
and $\text{depth}(M_{\mathfrak p}) = 0$. Since
$\text{depth}(M_{\mathfrak p}) = 0$, we see that
$\mathfrak p \in \text{Ass}(M)$ by the two Lemmas
\ref{lemma-associated-primes-localize} and \ref{lemma-ideal-nonzerodivisor}.
Since $\dim(\text{Supp}(M_{\mathfrak p})) \geq 1$, there
is a prime $\mathfrak q \in \text{Supp}(M)$ with
$\mathfrak q \subset \mathfrak p$, $\mathfrak q \not = \mathfrak p$.
We can take such a $\mathfrak q$ that is minimal in
$\text{Supp}(M)$. Then by
Proposition \ref{proposition-minimal-primes-associated-primes}
we have $\mathfrak q \in \text{Ass}(M)$ and hence
$\mathfrak p$ is an embedded associated prime.
\end{proof}

\begin{lemma}
\label{lemma-criterion-reduced}
\begin{slogan}
Reduced equals R0 plus S1.
\end{slogan}
Let $R$ be a Noetherian ring.
The following are equivalent:
\begin{enumerate}
\item $R$ is reduced, and
\item $R$ has properties $(R_0)$ and $(S_1)$.
\end{enumerate}
\end{lemma}

\begin{proof}
Suppose that $R$ is reduced. Then $R_{\mathfrak p}$ is a field for
every minimal prime $\mathfrak p$ of $R$, according to
Lemma \ref{lemma-minimal-prime-reduced-ring}. Hence we have $(R_0)$.
Let $\mathfrak p$ be a prime of height $\geq 1$. Then $A = R_{\mathfrak p}$
is a reduced local ring of dimension $\geq 1$. Hence its maximal
ideal $\mathfrak m$ is not an associated prime
since this would mean there exists an $x \in \mathfrak m$
with annihilator $\mathfrak m$ so $x^2 = 0$. Hence the depth of
$A = R_{\mathfrak p}$ is at least one, by Lemma \ref{lemma-ass-zero-divisors}.
This shows that $(S_1)$ holds.

\medskip\noindent
Conversely, assume that $R$ satisfies $(R_0)$ and $(S_1)$.
If $\mathfrak p$ is a minimal prime of $R$, then
$R_{\mathfrak p}$ is a field by $(R_0)$, and hence is reduced.
If $\mathfrak p$ is not minimal, then we see that $R_{\mathfrak p}$
has depth $\geq 1$ by $(S_1)$ and we conclude there exists an element
$t \in \mathfrak pR_{\mathfrak p}$ such that
$R_{\mathfrak p} \to R_{\mathfrak p}[1/t]$ is injective.
Now $R_\mathfrak p[1/t]$ is contained in the product
of its localizations at prime ideals, see
Lemma \ref{lemma-characterize-zero-local}.
This implies that $R_{\mathfrak p}$ is a subring of a product of
localizations of $R$ at $\mathfrak p \supset \mathfrak q$ with
$t \not \in \mathfrak q$. Since these primes have smaller height
by induction on the height we conclude that $R$ is reduced.
\end{proof}

\begin{lemma}[Serre's criterion for normality]
\label{lemma-criterion-normal}
\begin{reference}
\cite[IV, Theorem 5.8.6]{EGA}
\end{reference}
\begin{slogan}
Normal equals R1 plus S2.
\end{slogan}
Let $R$ be a Noetherian ring.
The following are equivalent:
\begin{enumerate}
\item $R$ is a normal ring, and
\item $R$ has properties $(R_1)$ and $(S_2)$.
\end{enumerate}
\end{lemma}

\begin{proof}
Proof of (1) $\Rightarrow$ (2). Assume $R$ is normal, i.e., all
localizations $R_{\mathfrak p}$ at primes are normal domains.
In particular we see that $R$ has $(R_0)$ and $(S_1)$ by
Lemma \ref{lemma-criterion-reduced}. Hence it suffices to show
that a local Noetherian normal domain $R$ of dimension $d$ has
depth $\geq \min(2, d)$ and is regular if $d = 1$. The assertion
if $d = 1$ follows from Lemma \ref{lemma-characterize-dvr}.

\medskip\noindent
Let $R$ be a local Noetherian normal domain with maximal ideal
$\mathfrak m$ and dimension $d \geq 2$. Apply
Lemma \ref{lemma-hart-serre-loc-thm} to $R$.
It is clear that $R$ does not fall into cases (1) or (2)
of the lemma.
Let $R \to R'$ as in (4) of the lemma.
Since $R$ is a domain we have $R \subset R'$. Since $\mathfrak m$
is not an associated prime of $R'$ there exists an $x \in \mathfrak m$
which is a nonzerodivisor on $R'$. Then $R_x = R'_x$ so
$R$ and $R'$ are domains with the same fraction field. But
finiteness of $R \subset R'$ implies every element of $R'$ is integral
over $R$ (Lemma \ref{lemma-finite-is-integral})
and we conclude that $R = R'$ as $R$ is normal.
This means (4) does not happen. Thus we get the remaining possibility
(3), i.e., $\text{depth}(R) \geq 2$ as desired.

\medskip\noindent
Proof of (2) $\Rightarrow$ (1). Assume $R$ satisfies $(R_1)$ and $(S_2)$.
By Lemma \ref{lemma-criterion-reduced} we conclude that $R$ is
reduced. Hence it suffices to show that if $R$ is a reduced local
Noetherian ring of dimension $d$ satisfying $(S_2)$ and $(R_1)$
then $R$ is a normal domain. If $d = 0$, the result is clear.
If $d = 1$, then the result follows from Lemma \ref{lemma-characterize-dvr}.

\medskip\noindent
Let $R$ be a reduced local Noetherian ring with maximal ideal
$\mathfrak m$ and dimension $d \geq 2$ which satisfies $(R_1)$ and
$(S_2)$. By Lemma \ref{lemma-characterize-reduced-ring-normal}
it suffices to show that $R$ is integrally closed in its
total ring of fractions $Q(R)$. Pick $x \in Q(R)$ which is integral
over $R$. Then $R' = R[x]$ is a finite ring extension of $R$
(Lemma \ref{lemma-characterize-finite-in-terms-of-integral}).
Because $\dim(R_\mathfrak p) < d$ for
every nonmaximal prime $\mathfrak p \subset R$
we have $R_\mathfrak p = R'_\mathfrak p$ by induction.
Hence the support of $R'/R$ is $\{\mathfrak m\}$.
It follows that $R'/R$ is annihilated by a power of $\mathfrak m$
(Lemma \ref{lemma-Noetherian-power-ideal-kills-module}).
By Lemma \ref{lemma-hart-serre-loc-thm} this
contradicts the assumption that the depth of $R$ is $\geq 2 = \min(2, d)$
and the proof is complete.
\end{proof}

\begin{lemma}
\label{lemma-regular-normal}
A regular ring is normal.
\end{lemma}

\begin{proof}
Let $R$ be a regular ring. By
Lemma \ref{lemma-criterion-normal}
it suffices to prove that $R$ is $(R_1)$ and $(S_2)$.
As a regular local ring is Cohen-Macaulay, see
Lemma \ref{lemma-regular-ring-CM},
it is clear that $R$ is $(S_2)$.
Property $(R_1)$ is immediate.
\end{proof}

\begin{lemma}
\label{lemma-normal-domain-intersection-localizations-height-1}
Let $R$ be a Noetherian normal domain with fraction field $K$. Then
\begin{enumerate}
\item for any nonzero $a \in R$ the quotient $R/aR$ has no embedded primes,
and all its associated primes have height $1$
\item
$$
R = \bigcap\nolimits_{\text{height}(\mathfrak p) = 1} R_{\mathfrak p}
$$
\item For any nonzero $x \in K$ the quotient $R/(R \cap xR)$
has no embedded primes, and all its associates primes have height $1$.
\end{enumerate}
\end{lemma}

\begin{proof}
By Lemma \ref{lemma-criterion-normal} we see that $R$ has $(S_2)$.
Hence for any nonzero element $a \in R$ we see that $R/aR$ has $(S_1)$
(use Lemma \ref{lemma-depth-in-ses} for example)
Hence $R/aR$ has no embedded primes
(Lemma \ref{lemma-criterion-no-embedded-primes}).
We conclude the associated primes of $R/aR$ are exactly
the minimal primes $\mathfrak p$ over $(a)$, which have height $1$
as $a$ is not zero (Lemma \ref{lemma-minimal-over-1}). This proves (1).

\medskip\noindent
Thus, given $b \in R$ we have $b \in aR$ if and only if
$b \in aR_{\mathfrak p}$ for every minimal prime $\mathfrak p$
over $(a)$ (see Lemma \ref{lemma-zero-at-ass-zero}).
These primes all have height $1$ as seen above so
$b/a \in R$ if and only if $b/a \in R_{\mathfrak p}$ for all
height 1 primes. Hence (2) holds.

\medskip\noindent
For (3) write $x = a/b$. Let $\mathfrak p_1, \ldots, \mathfrak p_r$
be the minimal primes over $(ab)$. These all have height 1 by the above.
Then we see that
$R \cap xR = \bigcap_{i = 1, \ldots, r} (R \cap xR_{\mathfrak p_i})$
by part (2) of the lemma. Hence $R/(R \cap xR)$ is a submodule of
$\bigoplus R/(R \cap xR_{\mathfrak p_i})$.
As $R_{\mathfrak p_i}$ is a discrete valuation ring (by property $(R_1)$
for the Noetherian normal domain $R$, see Lemma \ref{lemma-criterion-normal})
we have $xR_{\mathfrak p_i} = \mathfrak p_i^{e_i}R_{\mathfrak p_i}$
for some $e_i \in \mathbf{Z}$. Hence the direct sum is equal
to $\bigoplus_{e_i > 0} R/\mathfrak p_i^{(e_i)}$, see
Definition \ref{definition-symbolic-power}.
By Lemma \ref{lemma-symbolic-power-associated}
the only associated prime of the module
$R/\mathfrak p^{(n)}$ is $\mathfrak p$. Hence the set of associate primes
of $R/(R \cap xR)$ is a subset of $\{\mathfrak p_i\}$ and there are
no inclusion relations among them. This proves (3).
\end{proof}















